% !TeX root = ../../../main.tex
\subsubsection{Results: Adverse Event Time Prediction}
\label{ch03:sec:04_results_survival}

We begin by looking at the prediction performance of the survival models from Section~\ref{ch03:sec:methods_survival}.
Tables~\ref{ch03:tab:survival_3_6} and~\ref{ch03:tab:survival_9_12} report the time-dependent Concordance Index (\(C^{td}\)), Integrated Brier Score (IBS), and time-dependent Area Under the Curve (\(\text{AUC}^{td}\)) for each survival model--Cox proportional hazards (CoxPH), Random Survival Forest (RSF), DeepSurv, and DeepHit--across the four snapshot horizons. 

Across different prediction time snapshots, RSF consistently outperforms other models in terms of \(C^{td}\) and \(\text{AUC}^{td}\) while achieving lower IBS values, suggesting superior calibration and predictive reliability.
CoxPH and DeepSurv exhibit similar performance trends, with CoxPH maintaining slightly higher discriminative ability at later time points, whereas DeepSurv performs better at shorter horizons.

DeepHit underperforms across all time points, with lower \(C^{td}\) and \(\text{AUC}^{td}\) values and higher IBS scores, indicating weaker discrimination and calibration.
At 3 months, its \(C^{td}\) is 0.570 and \(\text{AUC}^{td}\) is 0.588, the lowest among all models, and this trend persists at later horizons, with \(C^{td}\) dropping to 0.540 by 12 months.

The superior performance of RSF suggests its strength in capturing nonlinear relationships and complex interactions among risk scores, demographic features, and adherence patterns.
In contrast, CoxPH is interpretable but limited by its linearity and proportional hazards assumptions, reducing its flexibility in capturing complex dependencies.
DeepSurv and DeepHit exhibit larger variability performance metrics, possibly due to hyperparameter sensitivity.

